\documentclass[12pt,a4paper]{report}
\usepackage[utf8]{inputenc}
\usepackage[T1]{fontenc}
\usepackage[margin=1in]{geometry}
\usepackage{setspace}
\usepackage{amsmath, amssymb, amsthm}
\usepackage{mathtools}
\usepackage{booktabs}
\usepackage{array}
\usepackage{longtable}
\usepackage{graphicx}
\usepackage{float}
\usepackage{caption}
\usepackage{subcaption}
\usepackage{pgfplots}
\pgfplotsset{compat=1.18}
\usepackage{tikz}
\usepackage{pgfplotstable}
\usetikzlibrary{patterns}
\usepgfplotslibrary{fillbetween}
\usepackage{xcolor}
\usepackage{mdframed}
\usepackage{enumitem}
\usepackage{multirow}
\usepackage{fancyhdr}
\usepackage{titlesec}
\usepackage{tocloft}
\usepackage{algorithm}
\usepackage{algpseudocode}
\usepackage{listings}
% natbib not used -- manual bibliography
\usepackage{hyperref}
\hypersetup{
	colorlinks=true,
	linkcolor=blue!60!black,
	citecolor=green!50!black,
	urlcolor=blue!70!black
}

% ── Page Style ───────────────────────────────────────────────
% ── Page Style ───────────────────────────────────────────────
\fancypagestyle{plain}{
	\fancyhf{}
	\fancyfoot[C]{\thepage}
	\renewcommand{\headrulewidth}{0pt}
}
\doublespacing

% ── Theorem Environments ─────────────────────────────────────
\theoremstyle{definition}
\newtheorem{definition}{Definition}[chapter]
\newtheorem{example}{Example}[chapter]

\theoremstyle{plain}
\newtheorem{theorem}{Theorem}[chapter]
\newtheorem{lemma}[theorem]{Lemma}
\newtheorem{proposition}[theorem]{Proposition}
\newtheorem{corollary}[theorem]{Corollary}

\theoremstyle{remark}
\newtheorem{remark}{Remark}[chapter]

% ── Highlighted Boxes ────────────────────────────────────────
\newmdenv[
backgroundcolor=blue!6,
linecolor=blue!40,
linewidth=1pt,
roundcorner=5pt,
innerleftmargin=10pt,
innerrightmargin=10pt,
innertopmargin=8pt,
innerbottommargin=8pt
]{keybox}

\newmdenv[
backgroundcolor=green!6,
linecolor=green!50!black,
linewidth=1pt,
roundcorner=5pt,
innerleftmargin=10pt,
innerrightmargin=10pt,
innertopmargin=8pt,
innerbottommargin=8pt
]{resultbox}

% ── Python Code Style ────────────────────────────────────────
\lstdefinestyle{pythonstyle}{
	language=Python,
	basicstyle=\ttfamily\small,
	keywordstyle=\color{blue!70!black}\bfseries,
	stringstyle=\color{green!50!black},
	commentstyle=\color{gray}\itshape,
	numberstyle=\tiny\color{gray},
	numbers=left,
	stepnumber=1,
	numbersep=8pt,
	backgroundcolor=\color{gray!8},
	frame=single,
	rulecolor=\color{gray!40},
	breaklines=true,
	breakatwhitespace=false,
	showstringspaces=false,
	tabsize=4,
	captionpos=b,
}

% ── Chapter Title Formatting ───────────────────────────────
\titleformat{\chapter}[display]
{\normalfont\Huge\bfseries\color{blue!60!black}\centering}
{CHAPTER \thechapter}
{20pt}
{\Huge}
\titlespacing*{\chapter}{0pt}{40pt}{40pt}

% ── Custom Commands ──────────────────────────────────────────
\newcommand{\lagr}{\mathcal{L}}
\newcommand{\R}{\mathbb{R}}
\newcommand{\taustar}{\tau^{*}}
\newcommand{\partialL}[1]{\dfrac{\partial \lagr}{\partial #1}}

% ============================================================


\begin{document}
	\section{Complete Optimization Problem}
	
	The complete constrained optimization problem is given by:
	
	\begin{equation}
		\begin{aligned}
			& \underset{\tau}{\text{maximize}} & & g(\tau) = \alpha + \beta\tau - \gamma\tau^2 \\
			& \text{subject to} & & \tau \geq \tau_{min} \\
			& & & \tau \leq \tau_{max} \\
			& & & \tau \in (0, 1) \\
			& & & \alpha > 0, \; \beta > 0, \; \gamma > 0
		\end{aligned}
		\label{eq:full_problem}
	\end{equation}
	
	\section{Optimization Method}
	\label{sec:kkt_solution}
	
	\subsubsection{The Karush-Kuhn-Tucker(KKT) Conditions}
	The following are step by step procedures of applying the model to the KKT optimization conditions:
	
	\subsubsection{Step 1: Re-write Constraints in Standard Form}
	
	For the KKT conditions to be applied, every constraint must be 
	written in the form $g_i(\tau) \leq 0$. We re-write the two 
	constraints of problem~\eqref{eq:full_problem} as follows:
	
	\begin{align}
		g_1(\tau) &= \tau_{min} - \tau \leq 0 \label{eq:g1}\\
		g_2(\tau) &= \tau - \tau_{max} \leq 0 \label{eq:g2}
	\end{align}
	
	\subsubsection{Step 2: Construct the Lagrangian}
	
	
	Since we are maximizing, the Lagrangian 
	is formed by:
	\\
	
	\begin{equation}
		\lagr(\tau, \mu_1, \mu_2) = g(\tau) - \mu_1 g_1(\tau) 
		- \mu_2 g_2(\tau)
	\end{equation}
	
	Substituting $g(\tau)$, $g_1(\tau)$, and $g_2(\tau)$, into the Lagrangian we have:
	\begin{equation}
		\lagr(\tau, \mu_1, \mu_2) = \alpha + \beta\tau 
		- \gamma\tau^2 - \mu_1(\tau_{min} - \tau) 
		- \mu_2(\tau - \tau_{max})
		\label{eq:lagrangian}
	\end{equation}
	
	where $\mu_1 \geq 0$ and $\mu_2 \geq 0$ are the KKT multipliers 
	for constraints $g_1$ and $g_2$ respectively. 
	
	\subsubsection{Step 3: Apply the Four KKT Conditions}
	
	For $\tau^*$ to be the optimal solution of 
	problem~\eqref{eq:full_problem}, the following four conditions 
	must all be satisfied simultaneously.
	
	\paragraph{Stationarity.}
	
	We differentiate the Lagrangian with respect to $\tau$ and set 
	it equal to zero:
	
	\begin{equation}
		\frac{\partial \lagr}{\partial \tau} = \beta - 2\gamma\tau 
		+ \mu_1 - \mu_2 = 0
		\label{eq:stationarity}
	\end{equation}
	
	This equation links the optimal tax rate $\taustar$ directly 
	to the two multipliers $\mu_1$ and $\mu_2$.
	
	\paragraph{Primal Feasibility.}
	
	The optimal solution must satisfy all the constraints:
	
	\begin{align}
		\tau_{min} - \tau &\leq 0 \quad \Longleftrightarrow \quad 
		\tau \geq \tau_{min} \label{eq:pf1}\\
		\tau - \tau_{max} &\leq 0 \quad \Longleftrightarrow \quad 
		\tau \leq \tau_{max} \label{eq:pf2}
	\end{align}
	
	\paragraph{Dual Feasibility.}
	
	The KKT multipliers must both be non-negative:
	
	\begin{equation}
		\mu_1 \geq 0, \qquad \mu_2 \geq 0
		\label{eq:df}
	\end{equation}
	
	\paragraph{Complementary Slackness.}
	
	For each constraint, either the constraint is active or its 
	multiplier is zero:
	
	\begin{align}
		\mu_1(\tau_{min} - \tau) &= 0 \label{eq:cs1}\\
		\mu_2(\tau - \tau_{max}) &= 0 \label{eq:cs2}
	\end{align}
	
	\subsubsection{Step 4: For Case Analysis}
	
	From the complementary slackness 
	conditions~\eqref{eq:cs1} and~\eqref{eq:cs2}, there are 
	three(3) important cases to consider:
	
	\subsubsection*{Case 1: Interior Solution 
		($\mu_1 = 0$, $\mu_2 = 0$, both constraints inactive)}
	
	If neither constraint is binding, both multipliers are zero. 
	Substituting $\mu_1 = \mu_2 = 0$ into the stationarity 
	condition~\eqref{eq:stationarity}:
	
	\begin{equation}
		\beta - 2\gamma\tau + 0 - 0 = 0
	\end{equation}
	
	\begin{equation}
		\beta - 2\gamma\tau = 0
	\end{equation}
	
	It implies:
	\begin{equation}
		\taustar = \frac{\beta}{2\gamma}
		\label{eq:interior}
	\end{equation}
	
	\begin{keybox}
		\textbf{Optimal Tax Rate for Case 1:}
		\[
		\taustar = \frac{\beta}{2\gamma}
		\]
		This is the tax rate at which the marginal benefit of 
		taxation (captured by $\beta$) exactly equals the marginal 
		cost of taxation (captured by $2\gamma\taustar$). This 
		solution is valid when $\tau_{min} < \taustar < \tau_{max}$.
	\end{keybox}
	
	We verify whether it is a maximum by checking the second derivative 
	of $g(\tau)$, so it implies:
	\[
	g''(\tau) = -2\gamma < 0 \quad \text{since } \gamma > 0 
	\quad \checkmark
	\]
	The negative second derivative confirms that $\taustar$ is 
	indeed a maximum.
	
	\subsubsection*{Case 2: Lower Bound Active 
		($\tau = \tau_{min}$, $\mu_2 = 0$)}
	
	If the unconstrained optimum lies below $\tau_{min}$, the 
	lower bound constraint becomes binding, meaning 
	$\tau^* = \tau_{min}$. Substituting $\tau = \tau_{min}$ 
	and $\mu_2 = 0$ into the stationarity 
	condition~\eqref{eq:stationarity}:
	
	\begin{equation}
		\beta - 2\gamma\tau_{min} + \mu_1 - 0 = 0
	\end{equation}
	
	\begin{equation}
		\mu_1 = 2\gamma\tau_{min} - \beta
		\label{eq:case2_mu1}
	\end{equation}
	
	For this case to be valid, dual feasibility requires 
	$\mu_1 \geq 0$:
	
	\[
	2\gamma\tau_{min} - \beta \geq 0 \implies 
	\frac{\beta}{2\gamma} \leq \tau_{min}
	\]
	
	This means the unconstrained optimum falls below the minimum 
	allowed tax rate. The lower bound constraint is therefore 
	active, and the constrained optimum is 
	$\taustar = \tau_{min}$.
	
	\subsubsection*{Case 3: Upper Bound Active 
		($\tau = \tau_{max}$, $\mu_1 = 0$)}
	
	If the unconstrained optimum exceeds $\tau_{max}$, the upper 
	bound constraint becomes binding, meaning 
	$\tau^* = \tau_{max}$. Substituting $\tau = \tau_{max}$ 
	and $\mu_1 = 0$ into the stationarity 
	condition~\eqref{eq:stationarity}:
	
	\begin{equation}
		\beta - 2\gamma\tau_{max} + 0 - \mu_2 = 0
	\end{equation}
	
	\begin{equation}
		\mu_2 = \beta - 2\gamma\tau_{max}
		\label{eq:case3_mu2}
	\end{equation}
	
	For this case to be valid, dual feasibility requires 
	$\mu_2 \geq 0$:
	
	\[
	\beta - 2\gamma\tau_{max} \geq 0 \implies 
	\frac{\beta}{2\gamma} \geq \tau_{max}
	\]
	
	This means the unconstrained optimum exceeds the maximum 
	allowed tax rate. The upper bound constraint is therefore 
	active, and the constrained optimum is 
	$\taustar = \tau_{max}$.
	
	\subsubsection*{Summary of All Three Cases}
	
	\begin{equation}
		\taustar =
		\begin{cases}
			\tau_{min} & \text{if } \dfrac{\beta}{2\gamma} 
			\leq \tau_{min} \\[10pt]
			\dfrac{\beta}{2\gamma} & \text{if } \tau_{min} < 
			\dfrac{\beta}{2\gamma} < \tau_{max} \\[10pt]
			\tau_{max} & \text{if } \dfrac{\beta}{2\gamma} 
			\geq \tau_{max}
		\end{cases}
		\label{eq:cases_summary}
	\end{equation}
	\subsubsection{Sufficiency of the KKT Conditions}
	
	\subsubsection{Step 5: Sufficiency Verification}
	Having identified the candidate optimal solution $\taustar$ 
	from the KKT conditions, we now confirm that it is indeed 
	a global maximum and not just any stationary point.
	
	We check the second derivative of the objective function 
	$g(\tau) = \alpha + \beta\tau - \gamma\tau^2$:
	
	\[
	g''(\tau) = -2\gamma
	\]
	
	Since $\gamma > 0$ by assumption, we have:
	
	\[
	g''(\tau) = -2\gamma < 0 \quad \text{for all } \tau
	\]
	
	This means $g(\tau)$ is strictly concave. 
	For a strictly concave objective function with convex constraints, the KKT conditions are both necessary and sufficient for a global maximum. Therefore, the solution 
	$\taustar$ obtained from the KKT conditions is guaranteed to 
	be a unique global maximum of the optimization 
	problem~\eqref{eq:full_problem}.
	
\end{document}
